% This is the ADASS_template.tex LaTeX file, 24 Oct 2025.
% It is based on the ASP general author template file, but modified to reflect the specific
% requirements of the ADASS proceedings.

% This is the aspauthor.tex LaTeX file
% Copyright 2021, Astronomical Society of the Pacific Conference Series
% Revision:  29 November 2021

% To compile, at the command line positioned at this folder, type:
% pdflatex ADASS_template
% pdflatex ADASS_template
% This will create a file called ADASS_template.pdf

\documentclass[11pt,twoside]{article}

% Do NOT use ANY packages other than asp2023. 
\usepackage{asp2023}

\aspSuppressVolSlug
\resetcounters

% References must all use BibTeX entries in a .bibfile.
% References must be cited in the text using \citet{} or \citep{}.
% Do not use \cite{}.
% See ManuscriptInstructions.pdf for more details
\bibliographystyle{asp2021}

% The ``markboth'' line sets up the running heads for the paper.
% 1 author: "Surname"
% 2 authors: "Surname1 and Surname2"
% 3 authors: "Surname1, Surname2, and Surname3"
% >3 authors: "Surname1 et al."
% Replace ``Short Title'' with the actual paper title, shortened if necessary.
% Use mixed case type for the shortened title
% Ensure shortened title does not cause an overfull hbox LaTeX error
% See ASPmanual2010.pdf 2.1.4  and ManuscriptInstructions.pdf for more details
\markboth{Bryant et al.}{The FIFI-LS Pipeline at the SOFIA Data Center}

\begin{document}

\title{The FIFI-LS Pipeline at the SOFIA Data Center - Preserving the legacy of Far-Infrared Imaging Spectroscopy}

% Note the position of the comma between the author name and the 
% affiliation number.
% Authors surnames should come after first names or initials, eg John Smith, or J. Smith.
% Author names should be separated by commas.
% The final author should be preceded by "and".
% Affiliations should not be repeated across multiple \affil commands. If several
% authors share an affiliation this should be in a single \affil which can then
% be referenced for several author names. If only one affiliation, no footnotes are needed.
% See ManuscriptInstructions.pdf and ASP's manual2010.pdf 3.1.4 for more details
\author{Aaron Bryant, Jonas Fr\"{u}h, Benjamin Greiner, Christof Iserlohe, and Bastian Knieling}
\affil{SOFIA Data Center, Institut f\"{u}r Raumfahrtsysteme, Universit\"{a}t Stuttgart, Pfaffenwaldring 31, 70569 Stuttgart, Germany; \email{aaron.bryant@irs.uni-stuttgart.de}}

% This section is for ADS Processing.  There must be one line per author. paperauthor has 9 arguments.
\paperauthor{Aaron Bryant}{aaron.bryant@irs.uni-stuttgart.de}{0000-0002-9472-2391}{Universit\"{a}t Stuttgart}{Institut f\"{u}r Raumfahrtsysteme}{Stuttgart}{Baden-W\"{u}rttemberg}{70569}{Germany}
\paperauthor{Jonas Fr\"{u}h}{jonas.frueh@irs.uni-stuttgart.de}{0009-0006-4927-3861}{Universit\"{a}t Stuttgart}{Institut f\"{u}r Raumfahrtsysteme}{Stuttgart}{Baden-W\"{u}rttemberg}{70569}{Germany}
\paperauthor{Benjamin Greiner}{benjamin.greiner@irs.uni-stuttgart.de}{0009-0008-1244-7711}{Universit\"{a}t Stuttgart}{Institut f\"{u}r Raumfahrtsysteme}{Stuttgart}{Baden-W\"{u}rttemberg}{70569}{Germany}
\paperauthor{Christof Iserlohe}{christof.iserlohe@irs.uni-stuttgart.de}{0000-0003-4223-7439}{Universit\"{a}t Stuttgart}{Institut f\"{u}r Raumfahrtsysteme}{Stuttgart}{Baden-W\"{u}rttemberg}{70569}{Germany}
\paperauthor{Bastian Knieling}{bastian.knieling@irs.uni-stuttgart.de}{0009-0000-7038-0406}{Universit\"{a}t Stuttgart}{Institut f\"{u}r Raumfahrtsysteme}{Stuttgart}{Baden-W\"{u}rttemberg}{70569}{Germany}

% There should be one \aindex line (commented out) for each author. These are used to
% build up the author index for the Proceedings. The surname must come first, followed by
% initials. Note the use of ~ before each initial to control spacing.
% The \author entries (see above) have surname last. These \aindex entries have
% surname first.
% The Aindex.py command willl create them for you after you have constructed the \author
% The first entry should be the first author, for bold-facing the author index page-reference

%\aindex{Bryant,~A.}
%\aindex{Fr\"{u}h,~J.}
%\aindex{Greiner,~B.}
%\aindex{Iserlohe,~C.}
%\aindex{Knieling,~B.}


\begin{abstract}
The SOFIA Data Center (SDC) was founded in 2024 at the Institut f\"{u}r Raumfahrtsysteme,
Universit\"{a}t Stuttgart, Germany. Its goal is to preserve the entire scientific and
technical databank of the former airborne observatory SOFIA, enhance the quality of
its contents, and support astronomical research through free and accessible data
distribution. Data reprocessing is enabled by an updated version of the data reduction
pipeline package "Redux", incorporating the experience of over a decade of science
mission operations and data analysis. The rich scientific legacy of SOFIA's observations
will be consequently made available to the global community in the form of a
"Virtual Observatory" (VO) hosted by the University of Stuttgart's Data Repository,
adhering to the ``FAIR'' principles of data assets.

Among the science pipelines being actively maintained at the SDC is that of the
imaging spectrometer FIFI-LS, which operated on SOFIA from 2014-2022, delivering
a wealth of far-infrared spectral data. Much of this data remains unpublished,
and contains astrophysical diagnostics from wavelength ranges that will be out
of reach to any observatory for at least the next half-decade. As of mid-2026,
work \& improvements to the FIFI-LS pipeline are steadily converging, with focus
shifting to the reprocessing and subsequent archival of the full dataset on the
future VO.

This paper will overview the general functionality of the FIFI-LS pipeline, and
present several major data reduction enhancements implemented at the SDC. These
include a full flux recalibration, enhanced on-the-fly telluric correction, and
the automatic generation of science-ready data products. Additionally, the science
data reprocessing \& archival strategy will be outlined.
\end{abstract}

% These lines show examples of subject index entries. At this stage these have to commented
% out, and need to be on separate lines. Eventually, they will be automatically uncommented
% and used to generate entries in the Subject Index at the end of the Proceedings volume.
% Don't leave these in! - replace them with ones relevant to your paper.

% These lines show examples of ASCL index entries. At this stage these have to commented
% out, and need to be on separate lines. Eventually, they will be automatically uncommented
% and used to generate entries in the ASCL Index at the end of the Proceedings volume.
% The ascl.py command will scan your paper on possible code names.
% Don't leave these in! - replace them with ones relevant to your paper.


\section{Pipeline Functionality}
The FIFI-LS pipeline is part of the Redux pipeline package, a collection of python-based
data reduction pipelines for five of SOFIA's science instruments. For most of SOFIA's
operational life, Redux was maintained and distributed by the Universities Space Research Association (USRA).
After the end of the SOFIA mission, Redux was transferred to the SDC, where it is now being actively maintained and improved,
and will be used for a complete reprocessing of the entire SOFIA science data archive (see \citet{P604_adassxxxvi}).

The FIFI-LS pipeline processes raw channel-specific data into science-ready data products,
using a series of calibration and reduction steps, built in a modular manner into the Redux
reduction framework. Raw data represent readouts of each channel of the FIFI-LS detector
arrays in units of [V/s], written as a FITS binary table. From this, the pipeline
produces a 3D spatial-spectral cube of calibrated data, over a wavelength band
at a single sky source.

The individual pipeline steps are outlined below, grouped by category:


\begin{itemize}
\setlength{\itemindent}{-1.5em}
\setlength{\parskip}{2pt}
\setlength{\parindent}{-1.5em}
\item \textbf{Reformatting}\smallskip

\emph{Check Headers}: Reads raw data and verifies FITS headers

\emph{Split Grating \& Chop}: Consolidates raw data by chop phase \& grating position

\emph{Fit Ramps}: Fits voltage readout ramps to convert data to units [ADU/s]
\item \textbf{Demodulation}\smallskip

\emph{Subtract Chops}: Subtracts off chop data from on chop data to remove sky background

\emph{Combine Nods}: Combines nod positions to remove radiative offset
\item \textbf{Pixel Calibration}\smallskip

\emph{Lambda Calibrate}: Uses grating formulae to generate wavelength data in units [\textmu m]

\emph{Spatial Calibrate}: Uses optical distortion model to generate sky coordinate data

\emph{Apply Flat}: Applies spatial and spectral flat field corrections

\emph{Combine Scans}: Combines grating positions to reconstruct full spectral scan
\item \textbf{Signal Calibration}\smallskip

\emph{Telluric Correct}: Applies WV/Ozone-dependent atmospheric correction models

\emph{Flux Calibrate}: Uses flux response to convert data to physical units $[\mathrm{W/m^2/Hz}]$

\emph{Correct Wave Shift}: Corrects for barycentric and solar system velocity shifts
\item \textbf{Resampling}\smallskip

\emph{Resample}: Interpolates data onto regularly spaced 3D spatial-spectral grid

\emph{Linefit}: Attempts to fit superposition of spectral line and continuum emission

\emph{Make Spectral Map}: Generates quicklook map and spectrum for archive usage
\end{itemize}

\section{Algorithm Enhancements}
The SDC can benefit from over a decade of science mission operations, instrument characterisation and
data analysis, as well as the experience of instrument scientists and data engineers. This has
led to several major improvements to the FIFI-LS pipeline.\smallskip

\textbf{\emph{Enhanced Telluric Correction}}: Telluric correction is now performed using on-the-fly computation
of water vapour and ozone column densities, based on the atmospheric conditions at the time and location
of each observation (\citet{FIFI_Fischer_2021}, \citet{FIFI_Iserlohe_2021}, \citet{FIFI_Iserlohe_2022}).
Data is sourced from the European Centre for Medium-Range Weather Forecasts (ECMWF), and an example is shown in Figure \ref{P816_f1}.
These values are used to select from a pre-computed library of ATRAN\_SDC transmission models, which are
based on the original ATRAN code, but have been modified to incorporate several fixes and upgrades - see \citet{P809_adassxxxvi} for details.
Furthermore, the pipeline is now able to interpolate between the selected models to generate
a synthetic spectrum closely matching the actual atmospheric conditions.\smallskip

\textbf{\emph{Flux Recalibration}}: The flux calibration of FIFI-LS data has been improved by using a new set of
flux response profiles, derived from a re-analysis of the original Mars calibration data. Details of
this process were presented in \citet{PO-30_adassxxxv}, and the response curves are now live in the
current version of Redux. Further efforts are ongoing to improve scatter in the input raw data and filter
outlying datapoints, which should improve the systematic uncertainty of the response curves.\smallskip

\textbf{\emph{Spectral Line Fitting}}: The aforementioned linefit module is a new addition to the pipeline, which
allows for the automatic generation of science-ready products such as line fluxes, continuum fluxes, and
line velocities. The wavelength-dependent model \textit{M} consisting of a Gaussian line profile superimposed on a
linear continuum is defined as:

\begin{displaymath}
M(\lambda) = \frac{A}{\sqrt{2\pi\sigma^2}} \exp\left[-\frac{(\lambda-\mu)^2}{2\sigma^2}\right] + C
\end{displaymath}

where \textit{A} is the integrated line flux, \textit{\textmu} is the line central wavelength, $\sigma$ is the line width,
and \textit{C} is the continuum level. An example plot of this output is shown in Figure \ref{P816_f2}.
Default fit parameters should be robust enough for the majority of FIFI-LS science observations where a line
was detected, while hand-configured parameters will be supplied for several specific projects.\smallskip

\textbf{\emph{Optimised Data Products}}: The output of the pipeline is a "Level 4" product consisting
of a fully calibrated 3D spatial-spectral cube, which is often comprised of observations from multiple flights
or even observing campaigns. Ongoing work is being done to optimise these products to correctly reflect
the combined data, remove unnecessary extensions and keywords, and streamline the data structure for
a better user experience and eventual integration into the Virtual Observatory.\smallskip

\textbf{\emph{Documentation \& Tutorials}}: Equally important as the software itself is the accompanying user and developer
documentation, which is being actively maintained and published with each pipeline version release, accessible
through the SDC GitHub repository\footnote{\url{https://github.com/SOFIA-Data-Center/sofia_redux}}, in both web and PDF formats. Efforts were made to clarify the functionality of pipeline
algorithms, and to highlight potential caveats in the data reduction process. Future versions will
incorporate content from numerous existing archive handbooks, hence consolidating instrument
documentation and knowhow into a single source. Furthermore, a series of Jupyter notebooks are being developed as tutorials
to demonstrate the usage of the pipeline and the analysis of FIFI-LS data.

\section{Archival Strategy}
At the SDC, a complete reprocessing workflow has been developed, which can recreate >99\% of the original
IRSA-hosted science archive pipeline products. Raw and reprocessed data products will be stored
at the University of Stuttgart's data repository DaRUS. The next major work package is the ingestion
of the reprocessed metadata into the SDC's Virtual Observatory (VO), utilising the GAVO Data Center Helper Suite (DaCHS) framework.
This will enable reprocessed FIFI-LS data to be discoverable and accessible through common VO protocols, such as TAP and SIA.
After the conclusion of the SDC project, ownership of the VO is planned to be transferred to the German Center for Astrophysics (DZA).


\articlefigure[width=0.6\textwidth]{P816_f1.eps}{P816_f1}{Example plot of ECMWF water vapour data for a specific timestamp during a SOFIA science flight. The corresponding flight path is displayed in white.}

\articlefigure[width=0.45\textwidth]{P816_f2.eps}{P816_f2}{Example plot of the pipeline linefit product, taken towards the Milky Way Galactic Centre, at the wavelength of the [OI] 63 \textmu m fine structure line.}

%\clearpage % To force this stuff to happen by this point in the text, otherwise these will probably end up after the references.

\acknowledgements The SDC is funded by the DLR Space Administration under grant number FKZ 50OK2404

\bibliography{P816}  % For BibTex

% if we have space left, we might add a conference photograph here. Leave commented for now.
% \bookpartphoto[width=1.0\textwidth]{foobar.eps}{FooBar Photo (Photo: Any Photographer)}

\end{document}
